\section{Asset Families in Detail}
\label{app:families}

\tabref{tab:families} gives the guarantee surface in tabular form. This appendix
records the reasoning behind the three rows most likely to surprise a reader, since
the table has room for the conclusion but not for the argument.

\para{Why tools have no live alias.} A tool is a callable with a signature, and its
callers bind to a specific \code{(name, version)} pair. Introducing an alias would
mean a caller's tool could change signature underneath it between calls --- a
strictly worse failure mode than an explicit version bump, because it converts a
compile-time-visible incompatibility into a runtime one. Tools therefore expose
read-only family history and no live-target indirection. The corollary is that tool
remediation \emph{does} require a caller change, and we accept that cost rather than
introducing silent signature drift.

\para{Why skill rollback creates a new version.} A skill has a mutable current
record plus immutable snapshots. Rolling back restores a prior snapshot \emph{as a
new current version} rather than moving a pointer backward. The reason is that the
skill's identity is the current record, and rewinding it would erase the fact that
the intervening version ever existed --- which is exactly the history a governance
system exists to preserve. The observable consequence is that a skill's version
number is monotonic even under rollback, which differs from the prompt path's alias
semantics and is listed separately in \tabref{tab:families} for that reason.

\para{Why test sets have version counters and per-example validity intervals.} A
test set is evidence, and evidence must be interpretable at the time a score was
computed. A monolithic version counter alone would force a new version for every
example edit, making the history unusably granular. Per-example validity intervals
let the corpus evolve while keeping any historical score interpretable against the
corpus as it stood: given a score and its timestamp, the exact example set that
produced it is recoverable. This is the mechanism that makes \figref{fig:chain}'s
durable-residue claim true for measurements specifically.

\para{Why the agent family exists at all.} An agent in \sys is not an executable ---
agents are composed elsewhere (\secref{sec:background}). It is the anchor record
that verification items, refinement jobs, and approvals attach to, which gives those
records a stable owner and gives operators one lever (\code{enabled}) that the
workers read to pause and resume activity for a whole agent. Modelling it as a
family rather than as a foreign key keeps the authority model uniform: pausing an
agent is a scope-checked, audited mutation like any other.
